\chapter{Glosarium Indonesia--Inggris}
\enchap{Indonesian--English Glossary}

\begin{center}
	\small
	\renewcommand{\arraystretch}{1.15}
	\begin{longtable}{@{}p{0.45\textwidth} p{0.45\textwidth}@{}}
		\toprule
		\textbf{Indonesia} & \textbf{English} \\
		\midrule
		\endhead
		aksi integral & integral action \\
		aturan Bryson & Bryson's rule \\
		bentuk kanonik & canonical form \\
		bentuk posisi / kecepatan (PID) & position / velocity form \\
		bidang $z$ & $z$-plane \\
		daerah konvergensi & region of convergence \\
		defisit sudut & angle deficiency \\
		derau proses / pengukuran & process / measurement noise \\
		ekuivalen ZOH & ZOH equivalent \\
		emulasi & emulation \\
		estimasi keadaan & state estimation \\
		filter \textit{anti-alias} & anti-aliasing filter \\
		frekuensi Nyquist & Nyquist frequency \\
		fungsi transfer pulsa & pulse transfer function \\
		galat keadaan tunak & steady-state error \\
		identifikasi sistem & system identification \\
		indeks kinerja & performance index \\
		keteramatan & observability \\
		keterkendalian & controllability \\
		kestabilan BIBO & BIBO stability \\
		kuadrat terkecil & least squares \\
		kuantisasi & quantisation \\
		kuantisasi koefisien & coefficient quantisation \\
		kurva reaksi & process reaction curve \\
		kutub dominan & dominant poles \\
		lewatan maksimum & maximum overshoot \\
		lingkaran satuan & unit circle \\
		margin fasa / penguatan & phase / gain margin \\
		pembagian panjang & long division \\
		pecahan parsial & partial fractions \\
		pencacahan & sampling \\
		pencacahan patologis & pathological sampling \\
		penempatan kutub & pole placement \\
		pengamat & observer \\
		penguatan kritis & ultimate gain \\
		penguatan referensi & reference gain \\
		perancangan langsung & direct design \\
		persamaan beda & difference equation \\
		persamaan karakteristik & characteristic equation \\
		persamaan Riccati aljabar diskrit & discrete algebraic Riccati equation \\
		prinsip pemisahan & separation principle \\
		rekonstruksi & reconstruction \\
		respons frekuensi & frequency response \\
		riak antar-cacah & intersample ripple \\
		ruang keadaan & state space \\
		syarat sudut / magnitudo & angle / magnitude condition \\
		tempat kedudukan akar & root locus \\
		tipe sistem & system type \\
		titik tetap / mengambang & fixed point / floating point \\
		transformasi bilinear & bilinear transformation \\
		transformasi similaritas & similarity transformation \\
		transformasi Z & Z-transform \\
		tunda komputasi & computational delay \\
		umpan balik keadaan & state feedback \\
		waktu cacah & sampling period \\
		waktu turun & settling time \\
		\bottomrule
	\end{longtable}
\end{center}
